\section{Constitutional Availability}
\label{sec:constitutional}

A rigorous comparison of alternative systems must distinguish between systems
that are immediately available under current constitutional and statutory law
and systems that require either congressional action, constitutional amendment,
or both.  This distinction is not merely procedural: it determines which
systems can be evaluated as near-term reforms versus long-term structural changes.

\subsection{Immediately Available: Multi-Member Districts}

Multi-member districts are constitutionally available without amendment.
The relevant authority is the federal single-member district statute, codified
at 2 U.S.C. \S2c, which requires that each state electing more than one
Representative do so from single-member districts.  This statute was enacted
by Congress and can be repealed or amended by Congress.  Because it is a
statutory rather than constitutional requirement, no Article~V process is needed.

The Fair Representation Act (H.R.~3863, 117th Congress) would have amended
2 U.S.C. \S2c to permit states to use multi-member districts with ranked-choice
voting \citep{fair2021representation}.  Its repeated introduction without passage
reflects political rather than constitutional obstacles.

The Voting Rights Act implications of multi-member districts are more complex.
Historical MMDs were used in some states to dilute minority voting power, and
\emph{Thornburg v.\ Gingles} \citep{gingles1986} largely prohibited submerging
minority populations in multi-member district majorities.  However, the
three-to-five-member districts contemplated by the Fair Representation Act are
specifically designed to create minority representation opportunities, not to
dilute them.  Paper E.1 finds that our MMD algorithm produces a minority-
representation surplus relative to the VRA benchmark --- the opposite of the
dilution concern.

\subsection{Immediately Available: County-Based Representation}

County representation is more complex.  In states where Congress grants
state legislatures discretion over the method of electing Representatives
(as it does under Article~I, \S4), a state could theoretically adopt county-
based representation for its congressional delegation.  However, the one-person
one-vote requirement derived from \emph{Reynolds v.\ Sims} \citep{reynolds1964}
and applied to congressional elections in \emph{Wesberry v.\ Sanders}
\citep{wesberry1964} would prohibit fractional-vote county representation:
the population disparities between counties are too large to satisfy the
``as nearly as is practicable'' equal-population standard.

A constitutional amendment specifying that House seats be allocated to counties
and that counties with fractional allocations receive weighted votes would be
required for the strict county representation system described in E.2.
A weaker version --- restricting district-drawing to preserve county lines
where population equality permits --- is already common and does not require
amendment.

\subsection{Requires Constitutional Amendment: National Redistricting}

National redistricting without state lines is constitutionally unavailable.
Article~I, \S2 is unambiguous: ``Representatives shall be apportioned among
the several States according to their respective Numbers.''  This requires
not merely apportionment by state but the drawing of districts within states.
\emph{Wesberry v.\ Sanders} \citep{wesberry1964} confirmed that congressional
districts must be drawn within individual states.

The research value of E.3 is precisely that it is a constitutionally prohibited
counterfactual.  Quantifying the compactness cost of the federalism constraint
(approximately 12\% PP improvement forgone) gives a precise answer to a
question that is otherwise answered only impressionistically: how much does
the state-boundary requirement cost us in geometric terms?

\subsection{Available but Contested: Partisan-Fairness Allocation}

Partisan-fairness seat allocation (E.5) is constitutionally available in the
sense that Congress could, under Article~I, \S4, impose an efficiency-gap
criterion on state redistricting plans.  The Supreme Court has held that
federal courts cannot impose partisan symmetry standards \citep{rucho2019},
but this does not prevent Congress from doing so by statute.

The contested question is whether a congressional proportionality mandate would
itself violate the Constitution.  The argument against is that by dictating
electoral outcomes (seat allocation) rather than drawing procedures, such a
mandate would interfere with states' reserved powers.  This argument is weaker
than it appears: Congress has broad authority under Article~I, \S4 to
``make or alter'' state regulations for congressional elections, and the
Court has consistently upheld congressional intervention in redistricting.

\subsection{Immediately Available: Algorithmic Single-Member Redistricting (DIA)}

The District-Independent Algorithm, which draws compact single-member districts
using census geography and the METIS graph partitioner without partisan data,
is immediately available.  It is a redistricting method, not a structural
reform, and requires no statutory change at the federal level.  Individual
states adopting algorithmic redistricting by statute or initiative would need
only to incorporate the method into their state code alongside existing criteria
(population equality, contiguity, VRA compliance).

This is the key asymmetry between the DIA and the alternatives: all five
experimental systems require either congressional action or constitutional
amendment to implement at the national scale, whereas the DIA can be adopted
state by state within the existing legal framework.  Table~\ref{tab:constitutional}
summarizes the constitutional status of each system.\footnote{\emph{Moore v.\
Harper} (2023) confirmed that state courts may enforce state constitutional
redistricting requirements, which creates new state-court review venues for all
systems listed as ``immediately available'' in Table~\ref{tab:constitutional}.
Federal courts, however, remain closed to partisan gerrymandering claims
post-\emph{Rucho}.  The District-Independent Algorithm (DIA) is a model
redistricting method; it has not been formally enacted as a statutory standard
in any state, though several states use algorithmic criteria that are
substantively similar.}

\begin{table}[ht]
\centering
\caption{Constitutional Availability of Each System}
\label{tab:constitutional}
\begin{tabular}{lll}
\toprule
\textbf{System} & \textbf{Status} & \textbf{Authority Required} \\
\midrule
Baseline DIA (B.1)          & Available    & State statute or initiative \\
E.1 Multi-member            & Available    & Federal statute (repeal 2 U.S.C. \S2c) \\
E.2 County representation   & Constrained  & Constitutional amendment (equal pop.) \\
E.3 National redistricting  & Unavailable  & Constitutional amendment (Art.~I \S2) \\
E.4 Partisan similarity     & Available    & State or federal statute \\
E.5 Partisan-fairness alloc.& Available    & Federal statute (Art.~I \S4) \\
E.6 International           & N/A          & Applicable national law varies \\
\bottomrule
\end{tabular}
\end{table}
